%=============================================================================
% AGENT 7: SELF-ACCELERATION OF A PSI-BUBBLE DEVICE
% IN AN EXTERNAL GRAVITATIONAL FIELD
% Cross-Term Force Analysis with Three Amplification Mechanisms
%=============================================================================

\documentclass[11pt]{article}
\usepackage{amsmath,amssymb,physics,geometry,booktabs,tcolorbox,graphicx}
\geometry{margin=1in}

\newcommand{\psifield}{\psi}
\newcommand{\etac}{\eta_c}
\newcommand{\ka}{k_a}
\newcommand{\kG}{\kappa_G}

\title{Self-Acceleration of a $\psi$-Bubble Device\\
in an External Gravitational Field:\\
Cross-Term Force and Amplification Analysis}
\author{Agent 7 --- DFD $\psi$-Bubble Drive Analysis}
\date{March 2026}

\begin{document}
\maketitle

\begin{abstract}
We compute the gravitational cross-term force on a $\psi$-bubble device
arising from the interference of the external gravitational field
$\psi_{\rm grav}$ with the device-generated perturbation $\delta\psi_{\rm EM}$.
This is Channel~4 from the five-channel thrust analysis. The DFD
stress-energy tensor contains a bilinear cross-term
$T^{ij}_{\rm cross} \propto \partial^i\psi_{\rm grav}\,\partial^j\delta\psi_{\rm EM}$
that produces a net force $F_{\rm cross} = \ka (U_{\rm EM}/c^2)\,g\,f(\text{geometry})$
where $\ka = 3/(8\alpha) \approx 51.4$.
For a static 10\,T magnet in 1\,L volume, the weight anomaly is
$\Delta w/w \sim 2.3 \times 10^{-12}$.
We analyse three amplification mechanisms --- Q-factor resonance, parametric
amplification, and volume scaling --- and show that the combination of
a 5\,m radius superconducting shell at 50\,T with parametric gain of
$G_{\rm par} \sim 10^4$ yields $F \sim 30$\,N, which is the first
combination reaching useful thrust. However, the parametric channel
requires $|\lambda - 1| \neq 0$ (Appendix~R), and the accidental bound
$|\lambda - 1| \lesssim 3 \times 10^{-5}$ means the saturated amplitude
could be many orders of magnitude smaller. We map out the full parameter
space and identify the critical unknowns.
\end{abstract}

\tableofcontents

%=============================================================================
\section{Theoretical Foundation: The Cross-Term Force}\label{sec:foundation}
%=============================================================================

\subsection{Field Decomposition}

In the DFD framework, the total scalar field is:
\begin{equation}\label{eq:decomposition}
\psi_{\rm tot}(\mathbf{r},t) = \psi_{\rm grav}(\mathbf{r})
  + \delta\psi_{\rm EM}(\mathbf{r},t),
\end{equation}
where $\psi_{\rm grav}$ is the external gravitational field (Earth's, the Sun's,
etc.) and $\delta\psi_{\rm EM}$ is the perturbation generated by the device's
electromagnetic configuration through the EM--$\psi$ coupling.

For Earth's surface:
\begin{equation}
\psi_{\rm grav} \approx \frac{\Phi_N}{c^2} = -\frac{GM_\oplus}{Rc^2}
  \approx -6.95 \times 10^{-10},
\end{equation}
with gradient:
\begin{equation}\label{eq:grad-psi-grav}
\nabla\psi_{\rm grav} = \frac{\mathbf{g}}{c^2},
  \qquad |\nabla\psi_{\rm grav}| = \frac{g}{c^2}
  = \frac{9.8}{9 \times 10^{16}} = 1.09 \times 10^{-16}\,\text{m}^{-1}.
\end{equation}

\subsection{The DFD Stress-Energy Tensor}

The $\psi$-field carries an effective stress-energy tensor (derived from the
DFD Lagrangian):
\begin{equation}
T^{ij}_\psi = \frac{c^4}{8\pi G}\left[
  \partial^i\psi\,\partial^j\psi
  - \frac{1}{2}\delta^{ij}(\nabla\psi)^2
\right].
\end{equation}

Substituting $\psi = \psi_{\rm grav} + \delta\psi_{\rm EM}$ and extracting
the bilinear cross term:
\begin{equation}\label{eq:cross-stress}
\boxed{
T^{ij}_{\rm cross} = \frac{c^4}{8\pi G}\left[
  \partial^i\psi_{\rm grav}\,\partial^j\delta\psi_{\rm EM}
  + \partial^i\delta\psi_{\rm EM}\,\partial^j\psi_{\rm grav}
  - \delta^{ij}\,\nabla\psi_{\rm grav}\cdot\nabla\delta\psi_{\rm EM}
\right].
}
\end{equation}

\subsection{Noether Lock and the Open-System Loophole}

For a \textit{closed} system, the Noether momentum theorem gives
$dP^i/dt = 0$. But a device in Earth's gravitational field is an
\textit{open} system: the external field $\psi_{\rm grav}$ breaks
translational symmetry. The net force on the device is the surface
integral:
\begin{equation}\label{eq:surface-force}
F^i = -\oint_S T^{ij}_{\rm cross}\,dS_j,
\end{equation}
where $S$ is a surface enclosing the device. This is the gravitational
analogue of the Lorentz force on a dielectric body in an external
electric field.

\subsection{Evaluation for a Localised $\delta\psi_{\rm EM}$}

If $\delta\psi_{\rm EM}$ is localised within volume $V$ and $\psi_{\rm grav}$
varies slowly over this volume (valid for any lab-scale device in Earth's
field), we can pull $\nabla\psi_{\rm grav}$ outside the integral.
Using Gauss's theorem on (\ref{eq:surface-force}):
\begin{equation}\label{eq:volume-force}
F^i = \frac{c^4}{8\pi G}\int_V \left[
  \partial^j\partial^i\psi_{\rm grav}\,\delta\psi_{\rm EM}
  + \partial^i\psi_{\rm grav}\,\nabla^2\delta\psi_{\rm EM}
  + \ldots
\right] d^3r.
\end{equation}

For the uniform-gradient approximation ($\partial^j\partial^i\psi_{\rm grav}
\approx 0$ over the device volume), the dominant term is:
\begin{equation}\label{eq:force-simplified}
\boxed{
F_{\rm cross} = \frac{c^4}{8\pi G}\,\frac{g}{c^2}
  \int_V \nabla^2(\delta\psi_{\rm EM})\,d^3r
  = \frac{c^2 g}{8\pi G}\oint_S \nabla(\delta\psi_{\rm EM})\cdot d\mathbf{S}.
}
\end{equation}

Alternatively, using the DFD field equation for $\delta\psi_{\rm EM}$
sourced by EM energy:
\begin{equation}
\nabla^2(\delta\psi_{\rm EM}) = -\frac{8\pi G}{c^4}\,\ka\,u_{\rm EM},
\end{equation}
where $u_{\rm EM}$ is the EM energy density and
$\ka = 3/(8\alpha) \approx 51.4$ is the self-coupling coefficient
(Appendix~G, Theorem~\ref{thm:ka} of the Unified Theory). Substituting:
\begin{equation}\label{eq:force-ka}
\boxed{
F_{\rm cross} = -\ka\,\frac{g}{c^2}\int_V u_{\rm EM}\,d^3r
  = -\ka\,\frac{U_{\rm EM}}{c^2}\,g.
}
\end{equation}

\subsection{Baseline Result: Static Magnet}

For a 10\,T field in volume $V = 1$\,L $= 10^{-3}$\,m$^3$:
\begin{align}
U_{\rm EM} &= \frac{B^2}{2\mu_0}V
  = \frac{100}{2 \times 4\pi \times 10^{-7}} \times 10^{-3}
  = 39.8\,\text{kJ}, \\
F_{\rm cross} &= 51.4 \times \frac{39800}{9 \times 10^{16}} \times 9.8
  = 51.4 \times 4.42 \times 10^{-13} \times 9.8 \\
  &= 2.23 \times 10^{-10}\,\text{N}.
\end{align}

The weight of the magnet system (say $m_{\rm device} = 10$\,kg):
\begin{equation}
\frac{\Delta w}{w} = \frac{F_{\rm cross}}{m_{\rm device}\,g}
  = \frac{2.23 \times 10^{-10}}{98} = 2.3 \times 10^{-12}.
\end{equation}

This is far below measurement sensitivity. We now examine three
amplification mechanisms.


%=============================================================================
\section{Amplification Mechanism 1: Q-Factor Resonance}\label{sec:Q-amp}
%=============================================================================

\subsection{The Question}

A ferrite-superconductor resonant cavity with quality factor $Q > 10^9$
stores EM energy that is $Q$ times the energy dissipated per cycle.
If the EM field is driven at frequency $\omega$, the stored energy is:
\begin{equation}
U_{\rm stored} = Q \times \frac{P_{\rm input}}{\omega},
\end{equation}
where $P_{\rm input}$ is the input power maintaining steady state.

Does the cross-term force (\ref{eq:force-ka}) scale with the
\textit{stored} energy or the \textit{input} power?

\subsection{Analysis}

The force formula (\ref{eq:force-ka}) involves $U_{\rm EM} = \int u_{\rm EM}\,d^3r$,
which is the \textit{instantaneous total EM energy} in the cavity. In steady
state, this equals the stored energy $U_{\rm stored}$. Therefore:

\begin{equation}\label{eq:Q-force}
F_{\rm cross}^{(Q)} = \ka\,\frac{U_{\rm stored}}{c^2}\,g
  = \ka\,\frac{Q\,P_{\rm input}}{\omega\,c^2}\,g.
\end{equation}

\paragraph{Critical subtlety.}
However, this analysis is \textit{exactly} the static case rewritten.
The $Q$-factor does not provide additional amplification beyond increasing
the stored energy. If we simply build a stronger static magnet with the
same $U_{\rm EM}$, we get the same force. The $Q$-factor is merely an
efficient way to store large EM energy with modest input power.

\subsection{Numerical Evaluation}

For $Q = 10^9$, $P_{\rm input} = 13$\,kW, $\omega = 2\pi \times 10^8$\,Hz:
\begin{align}
U_{\rm stored} &= \frac{10^9 \times 1.3 \times 10^4}{6.28 \times 10^8}
  = 2.07 \times 10^4\,\text{J} = 20.7\,\text{kJ}, \\
F_{\rm cross}^{(Q)} &= 51.4 \times \frac{2.07 \times 10^4}{9 \times 10^{16}}
  \times 9.8 = 1.16 \times 10^{-10}\,\text{N}.
\end{align}

\begin{tcolorbox}[colback=red!5!white, colframe=red!50!black, title=Q-Factor Verdict]
\textbf{Q-factor resonance does not amplify the cross-term force
beyond the stored energy.} The force is $F = \ka(U_{\rm stored}/c^2)g$
regardless of whether the energy is stored via a static field or a
high-$Q$ oscillating cavity. For $U_{\rm stored} = 20$\,kJ, the force
is $\sim 10^{-10}$\,N.

The $Q$-factor is relevant only insofar as it allows large $U_{\rm stored}$
with small $P_{\rm input}$.
\end{tcolorbox}


%=============================================================================
\section{Amplification Mechanism 2: Parametric Resonance}\label{sec:parametric}
%=============================================================================

\subsection{The Mechanism}

From Appendix~R of the Unified Theory, the $\psi$-mode equation with
parametric driving is (Eq.~R.1):
\begin{equation}\label{eq:mode-eq}
\ddot{q} + 2\gamma_\psi\dot{q} + \Omega_\psi^2 q
  = \frac{(\lambda - 1)}{M_\psi}G_{\rm driven}
  + \alpha_{\rm par} U(t)\,q,
\end{equation}
where the stiffness modulation from $U(t) = U_0[1 + m\cos(2\omega t)]$
creates a Mathieu-type instability when $2\omega \simeq 2\Omega_\psi$.

The Mathieu parameter is:
\begin{equation}
h = (\lambda - 1)\frac{U_0\,H\,m}{M_\psi\Omega_\psi^2},
\end{equation}
and the parametric growth rate:
\begin{equation}\label{eq:growth-rate}
\Gamma = \frac{1}{2}h\,\Omega_\psi - \gamma_\psi.
\end{equation}

Above the instability threshold ($\Gamma > 0$), the $\psi$-mode amplitude
grows exponentially: $q(t) \propto e^{\Gamma t}$.

\subsection{Saturated Amplitude}

Exponential growth is ultimately limited by nonlinearity. In the DFD
framework, the leading nonlinear term in the $\psi$ field equation is the
self-interaction at order $\psi^2$, with coefficient scaling as $\ka$.
The saturated amplitude is estimated by balancing the parametric gain
against nonlinear detuning:
\begin{equation}\label{eq:saturated-q}
|q|_{\rm sat} \sim \sqrt{\frac{\Gamma}{\ka\,\Omega_\psi^2}}
  = \sqrt{\frac{h\Omega_\psi/2 - \gamma_\psi}{\ka\,\Omega_\psi^2}}.
\end{equation}

For $\Gamma \gg \gamma_\psi$ (well above threshold):
\begin{equation}
|q|_{\rm sat} \sim \sqrt{\frac{h}{2\ka\,\Omega_\psi}}.
\end{equation}

\subsection{How Parametric Resonance Amplifies the Cross-Term}

The key insight: in the parametric regime, the $\psi$-mode amplitude
$|q|_{\rm sat}$ can be \textit{much larger} than the directly driven
amplitude $|q|_{\rm driven}$. The ratio defines the parametric gain:
\begin{equation}\label{eq:par-gain}
G_{\rm par} \equiv \frac{|q|_{\rm sat}}{|q|_{\rm driven}}.
\end{equation}

The device-generated perturbation becomes:
\begin{equation}
\delta\psi_{\rm EM} = q(t)\,u(\mathbf{r}) \to |q|_{\rm sat}\,u(\mathbf{r}),
\end{equation}
and the cross-term force (\ref{eq:force-simplified}) scales with
$\delta\psi_{\rm EM}$, so:
\begin{equation}\label{eq:amplified-force}
\boxed{
F_{\rm cross}^{(\rm par)} = G_{\rm par} \times F_{\rm cross}^{(\rm static)}.
}
\end{equation}

\subsection{Estimating the Parametric Gain}

The directly driven $\delta\psi_{\rm EM}$ for a static magnet is:
\begin{equation}
\delta\psi_{\rm driven} \sim \ka\,\frac{U_{\rm EM}}{M_\psi c^2},
\end{equation}
where $M_\psi c^2 \sim c^4/(G\Omega_\psi^2 V_\psi)$ is the effective
mass-energy of the $\psi$ mode of volume $V_\psi$.

The saturated parametric amplitude:
\begin{equation}
\delta\psi_{\rm sat} \sim |q|_{\rm sat}\,u_{\rm max}
  \sim \sqrt{\frac{h}{2\ka\,\Omega_\psi}}.
\end{equation}

For the parametric gain:
\begin{equation}\label{eq:gain-estimate}
G_{\rm par} \sim \frac{\delta\psi_{\rm sat}}{\delta\psi_{\rm driven}}
  \sim \frac{1}{\delta\psi_{\rm driven}}
    \sqrt{\frac{(\lambda - 1) U_0 H m}{2\ka\,M_\psi\Omega_\psi^3}}.
\end{equation}

\subsection{The $\lambda$ Dependence Problem}

\textbf{This is the critical issue.} The parametric gain depends on
$(\lambda - 1)$:
\begin{itemize}
\item If $\lambda = 1$ exactly: $h = 0$, no parametric gain, $G_{\rm par} = 1$.
\item If $|\lambda - 1| = 3 \times 10^{-5}$ (accidental bound): marginal
  parametric gain depending on whether $\Gamma > 0$.
\item If $|\lambda - 1| \sim 10^{-14}$ (intentional search reach): extremely
  small gain, deep below threshold for any realistic cavity.
\end{itemize}

The instability threshold (Eq.~R.7 of the Unified Theory) is:
\begin{equation}
|\lambda - 1|_{\rm min} = \frac{2\gamma_\psi}{\Omega_\psi}
  \frac{M_\psi\Omega_\psi^2}{U_0 H m}.
\end{equation}

For the design-law version (Eq.~R.10):
\begin{equation}
|\lambda - 1|_{\rm min} = \frac{\pi\gamma_\psi}{c_s U_0 m}
  \frac{A_\psi^2}{\kappa_{\rm eff} A_{\rm cav,tot}}.
\end{equation}

\paragraph{Numerical evaluation for a large cavity.}
Using intentional-search parameters from Appendix~R Table~1:
$U_0 = 10^6$\,J, $m = 0.1$, $A_{\rm cav,tot} = 0.03$\,m$^2$,
$A_\psi = 0.27$\,m$^2$, $\gamma_\psi/\Omega_\psi = 10^{-3}$,
$c_s = c$:
\begin{equation}
|\lambda - 1|_{\rm min}
  = \frac{\pi \times 10^{-3} \times \Omega_\psi}{3 \times 10^8 \times 10^6
    \times 0.1}
  \frac{0.0729}{1 \times 0.03}
  \approx 10^{-14}.
\end{equation}

So the instability threshold is $|\lambda - 1| \sim 10^{-14}$, meaning
parametric amplification occurs only if $|\lambda - 1| > 10^{-14}$.
Given the accidental bound $|\lambda - 1| \lesssim 3 \times 10^{-5}$,
there is a \textit{window} of nine orders of magnitude where parametric
gain could operate.

\subsection{Parametric Gain as a Function of $|\lambda - 1|$}

If $|\lambda - 1|$ is at the accidental bound ($3 \times 10^{-5}$) and
the threshold is $10^{-14}$, then:
\begin{equation}
\frac{h}{h_{\rm min}} = \frac{|\lambda - 1|}{|\lambda - 1|_{\rm min}}
  = \frac{3 \times 10^{-5}}{10^{-14}} = 3 \times 10^9.
\end{equation}

The growth rate:
\begin{equation}
\Gamma = \frac{1}{2}h\,\Omega_\psi - \gamma_\psi
  = \gamma_\psi\left(\frac{h}{h_{\rm min}} - 1\right)
  \approx 3 \times 10^9 \times \gamma_\psi.
\end{equation}

This is enormously above threshold. The saturated amplitude from
(\ref{eq:saturated-q}):
\begin{equation}
|q|_{\rm sat} \sim \sqrt{\frac{3 \times 10^9 \times \gamma_\psi}
  {\ka\,\Omega_\psi^2}}.
\end{equation}

However, this must be compared with the \textit{linear} regime. The
directly driven amplitude (from the static force) is:
\begin{equation}
\delta\psi_{\rm driven}
  = \ka\,\frac{U_{\rm EM}}{c^4/(G\,\Omega_\psi^2\,V_\psi)}
  = \ka\,\frac{G\,U_{\rm EM}\,\Omega_\psi^2\,V_\psi}{c^4}.
\end{equation}

The parametric gain factor:
\begin{equation}
G_{\rm par} = \frac{|q|_{\rm sat}}{\delta\psi_{\rm driven}/u_{\rm max}}
  \sim \sqrt{\frac{h/h_{\rm min}}{\ka}} \times \frac{1}{\delta\psi_{\rm driven}}
  \times \sqrt{\frac{\gamma_\psi}{\Omega_\psi}}.
\end{equation}

A more direct estimate: the saturated $\delta\psi$ is limited by
nonlinearity at order:
\begin{equation}\label{eq:psi-sat-estimate}
|\delta\psi_{\rm sat}| \sim \frac{1}{\sqrt{\ka}}
  \sim \frac{1}{\sqrt{51.4}} \approx 0.14.
\end{equation}

This is the \textit{absolute maximum} saturated amplitude from parametric
instability --- the nonlinear self-coupling caps $\psi$ perturbations at
$\sim 1/\sqrt{\ka}$. In practice, the saturation is much more nuanced and
depends on mode coupling, but this provides an upper bound.

\subsection{Upper Bound on Parametric Force}

If $|\delta\psi|_{\rm sat} \sim 0.14$ (the nonlinear ceiling), the
cross-term force is:
\begin{align}
F_{\rm cross}^{(\rm par,max)}
  &= \frac{c^4}{8\pi G}\,\frac{g}{c^2}
     \int_V |\nabla^2(\delta\psi_{\rm sat})|\,d^3r \notag \\
  &\sim \frac{c^2\,g}{8\pi G}\,
     \frac{|\delta\psi_{\rm sat}|}{L_\psi^2}\,V_\psi,
\end{align}
where $L_\psi$ is the characteristic scale of the $\psi$ mode and
$V_\psi \sim L_\psi^3$. Thus:
\begin{equation}
F_{\rm cross}^{(\rm par,max)}
  \sim \frac{c^2\,g}{8\pi G}\,|\delta\psi_{\rm sat}|\,L_\psi.
\end{equation}

Numerically, for $L_\psi = 5$\,m (a large device):
\begin{equation}
F_{\rm cross}^{(\rm par,max)}
  = \frac{(3 \times 10^8)^2 \times 9.8}{8\pi \times 6.674 \times 10^{-11}}
    \times 0.14 \times 5
  = \frac{8.82 \times 10^{17}}{1.676 \times 10^{-9}} \times 0.7
  \approx 3.7 \times 10^{26}\,\text{N}.
\end{equation}

This absurdly large number indicates that the nonlinear saturation estimate
(\ref{eq:psi-sat-estimate}) is far too generous. The actual saturation
must account for the fact that $\delta\psi_{\rm EM}$ is sourced by the EM
field, and the EM field itself has finite energy. The correct bound comes
from energy conservation.

\subsection{Energy-Conservation Bound on Saturated Amplitude}

The total energy available for $\psi$-mode pumping is bounded by the
EM stored energy $U_{\rm EM}$. The energy in the $\psi$ mode is:
\begin{equation}
E_\psi = \frac{1}{2}M_\psi\Omega_\psi^2|q|^2
  = \frac{c^4}{16\pi G}\int_V (\nabla\delta\psi)^2\,d^3r
  \sim \frac{c^4}{16\pi G}\frac{|\delta\psi|^2}{L_\psi^2}V_\psi.
\end{equation}

Energy conservation requires $E_\psi \leq U_{\rm EM}$:
\begin{equation}
|\delta\psi|_{\rm max}^{(\rm energy)} \leq
  \sqrt{\frac{16\pi G\,U_{\rm EM}\,L_\psi^2}{c^4\,V_\psi}}
  = \sqrt{\frac{16\pi G\,U_{\rm EM}}{c^4\,L_\psi}}.
\end{equation}

For $U_{\rm EM} = 10^6$\,J (1\,MJ stored energy), $L_\psi = 5$\,m:
\begin{equation}
|\delta\psi|_{\rm max}^{(\rm energy)}
  = \sqrt{\frac{16\pi \times 6.674 \times 10^{-11} \times 10^6}
  {(3 \times 10^8)^4 \times 5}}
  = \sqrt{\frac{3.35 \times 10^{-4}}{4.05 \times 10^{34}}}
  = \sqrt{8.3 \times 10^{-39}}
  = 2.9 \times 10^{-20}.
\end{equation}

This is the fundamental limit: even if all EM energy is perfectly converted
to $\psi$-mode energy, the perturbation is $|\delta\psi| \sim 3 \times 10^{-20}$.

The corresponding cross-term force:
\begin{align}
F_{\rm cross}^{(\rm par,energy)}
  &\sim \frac{c^2\,g}{8\pi G}\,\frac{|\delta\psi|_{\rm max}}{L_\psi} \notag \\
  &= \frac{(9 \times 10^{16})(9.8)}{8\pi(6.674 \times 10^{-11})}
    \times \frac{2.9 \times 10^{-20}}{5} \notag \\
  &= \frac{8.82 \times 10^{17}}{1.676 \times 10^{-9}}
    \times 5.8 \times 10^{-21} \notag \\
  &= 5.26 \times 10^{26} \times 5.8 \times 10^{-21} \notag \\
  &= 3.1 \times 10^{6}\,\text{N}.
\end{align}

Wait --- this is still enormous. The issue is that $c^4/(8\pi G)$ is
a huge prefactor ($\sim 5 \times 10^{26}$\,N/m$^2$). Let us re-examine
the force formula more carefully.

\subsection{Correct Force Formula from First Principles}

The cross-term force (\ref{eq:force-ka}) derived via the field equation
gives:
\begin{equation}
F_{\rm cross} = \ka\,\frac{U_{\rm EM}}{c^2}\,g.
\end{equation}

This is the \textit{correct} expression for a \textit{static} EM configuration.
In the parametric regime, $\delta\psi_{\rm EM}$ is amplified beyond the
static value. The amplified force should be expressed as:
\begin{equation}
F_{\rm cross}^{(\rm amplified)}
  = \frac{c^2\,g}{8\pi G}\oint_S
    \nabla(\delta\psi_{\rm EM}^{(\rm sat)})\cdot d\mathbf{S},
\end{equation}
which, using $\nabla^2(\delta\psi) \sim |\delta\psi|/L_\psi^2$ for a
mode of scale $L_\psi$:
\begin{equation}
F_{\rm cross}^{(\rm amplified)}
  \sim \frac{c^2\,g}{8\pi G}\,\frac{|\delta\psi_{\rm sat}|}{L_\psi^2}\,V_\psi
  = \frac{c^2\,g}{8\pi G}\,|\delta\psi_{\rm sat}|\,L_\psi.
\end{equation}

But note: the \textit{static} $\delta\psi$ sourced by $U_{\rm EM}$ is:
\begin{equation}
\delta\psi_{\rm static} \sim \ka\,\frac{G\,U_{\rm EM}}{c^4\,L_\psi},
\end{equation}
and the static force is:
\begin{equation}
F_{\rm static} = \frac{c^2 g}{8\pi G}\,\delta\psi_{\rm static}\,L_\psi
  = \frac{c^2 g}{8\pi G}\,\ka\,\frac{G\,U_{\rm EM}}{c^4\,L_\psi}\,L_\psi
  = \frac{\ka\,g\,U_{\rm EM}}{8\pi c^2}.
\end{equation}

Hmm, this differs from (\ref{eq:force-ka}) by a factor of $8\pi$.
The discrepancy arises because (\ref{eq:force-ka}) used the
Poisson-equation form where $\nabla^2\psi = -(8\pi G/c^4)\ka\,u_{\rm EM}$,
and the integral $\int u_{\rm EM}\,dV = U_{\rm EM}$. Let us be precise.

Starting from:
\begin{equation}
F_{\rm cross} = \frac{c^4}{8\pi G}\,\frac{g}{c^2}
  \int_V \nabla^2(\delta\psi_{\rm EM})\,d^3r.
\end{equation}

Using $\nabla^2(\delta\psi) = -(8\pi G/c^4)\ka\,u_{\rm EM}$:
\begin{equation}
F_{\rm cross} = \frac{c^4}{8\pi G}\,\frac{g}{c^2}
  \times\left(-\frac{8\pi G}{c^4}\right)\ka\,U_{\rm EM}
  = -\ka\,\frac{g\,U_{\rm EM}}{c^2}.
\end{equation}

This confirms (\ref{eq:force-ka}). The parametric amplification ratio is then:
\begin{equation}
\frac{F_{\rm par}}{F_{\rm static}}
  = \frac{\int \nabla^2(\delta\psi_{\rm sat})\,d^3r}
    {\int \nabla^2(\delta\psi_{\rm static})\,d^3r}
  = \frac{|\delta\psi_{\rm sat}|}{|\delta\psi_{\rm static}|} \equiv G_{\rm par}.
\end{equation}

Now $\delta\psi_{\rm static}$ for $U_{\rm EM} = 10^6$\,J, $L_\psi = 5$\,m:
\begin{equation}
\delta\psi_{\rm static} = \ka\,\frac{G\,U_{\rm EM}}{c^4\,L_\psi}
  = 51.4 \times \frac{6.674 \times 10^{-11} \times 10^6}{8.1 \times 10^{33} \times 5}
  = 51.4 \times 1.65 \times 10^{-39}
  = 8.5 \times 10^{-38}.
\end{equation}

And $\delta\psi_{\rm sat}^{(\rm energy)} = 2.9 \times 10^{-20}$, so:
\begin{equation}
G_{\rm par}^{(\rm max)} = \frac{2.9 \times 10^{-20}}{8.5 \times 10^{-38}}
  = 3.4 \times 10^{17}.
\end{equation}

The maximum parametric force:
\begin{equation}
F_{\rm par}^{(\rm max)} = 3.4 \times 10^{17} \times F_{\rm static}
  = 3.4 \times 10^{17} \times \ka\,\frac{U_{\rm EM}}{c^2}\,g.
\end{equation}

For $U_{\rm EM} = 10^6$\,J:
\begin{equation}
F_{\rm static} = 51.4 \times \frac{10^6}{9 \times 10^{16}} \times 9.8
  = 5.6 \times 10^{-9}\,\text{N}.
\end{equation}

\begin{equation}
F_{\rm par}^{(\rm max)} = 3.4 \times 10^{17} \times 5.6 \times 10^{-9}
  = 1.9 \times 10^{9}\,\text{N}.
\end{equation}

But this assumes \textit{all} EM energy converts to $\psi$-mode energy,
which violates steady-state operation. In practice, the conversion
efficiency $\epsilon$ is bounded by:
\begin{equation}
\epsilon \leq \frac{2\gamma_\psi}{\Omega_\psi} \times |\lambda - 1| \times H
  \ll 1.
\end{equation}

For $|\lambda - 1| = 3 \times 10^{-5}$,
$\gamma_\psi/\Omega_\psi = 10^{-3}$, $H \sim 1$:
\begin{equation}
\epsilon \lesssim 6 \times 10^{-8}.
\end{equation}

The realistic saturated $\delta\psi$:
\begin{equation}
|\delta\psi_{\rm sat}| \sim \sqrt{\epsilon}\,|\delta\psi|_{\rm max}^{(\rm energy)}
  \sim \sqrt{6 \times 10^{-8}} \times 2.9 \times 10^{-20}
  = 2.4 \times 10^{-4} \times 2.9 \times 10^{-20}
  = 7.1 \times 10^{-24}.
\end{equation}

The realistic parametric gain:
\begin{equation}
G_{\rm par}^{(\rm realistic)}
  = \frac{7.1 \times 10^{-24}}{8.5 \times 10^{-38}}
  = 8.4 \times 10^{13}.
\end{equation}

And the force:
\begin{equation}
F_{\rm par}^{(\rm realistic)}
  = 8.4 \times 10^{13} \times 5.6 \times 10^{-9}
  = 4.7 \times 10^{5}\,\text{N} = 470\,\text{kN}.
\end{equation}

\begin{tcolorbox}[colback=yellow!5!white, colframe=yellow!50!black,
  title=Parametric Resonance: Enormous but Conditional]
If $|\lambda - 1| \sim 3 \times 10^{-5}$ and all other parameters are at
their intentional-search values, parametric resonance could in principle
amplify the cross-term force by $\sim 10^{13}$ relative to static, giving
forces of order $10^5$\,N.

\textbf{However:}
\begin{enumerate}
\item The value of $|\lambda - 1|$ is unknown. If $\lambda = 1$ exactly,
  $G_{\rm par} = 1$ and no amplification occurs.
\item The $\epsilon$ estimate above is rough. The actual coupling efficiency
  depends on mode overlap integrals that could be much smaller.
\item The nonlinear saturation physics is complex and may cap the amplitude
  far below our estimate.
\item Maintaining parametric resonance requires exquisite frequency matching
  ($2\omega = 2\Omega_\psi$ to within $\gamma_\psi$).
\end{enumerate}
\end{tcolorbox}

%=============================================================================
\section{Amplification Mechanism 3: Volume Scaling}\label{sec:volume}
%=============================================================================

\subsection{Scaling Law}

The cross-term force (\ref{eq:force-ka}) scales linearly with
$U_{\rm EM} = u_{\rm EM} \times V$:
\begin{equation}
F_{\rm cross} = \ka\,\frac{u_{\rm EM}\,V}{c^2}\,g.
\end{equation}

For fixed field strength $B$ (fixed $u_{\rm EM} = B^2/(2\mu_0)$),
the force scales as $V$:
\begin{equation}
\frac{F(V_2)}{F(V_1)} = \frac{V_2}{V_1}.
\end{equation}

\subsection{Tabletop to Large Shell}

\begin{center}
\begin{tabular}{lcccc}
\toprule
Configuration & $B$ (T) & $V$ (m$^3$) & $U_{\rm EM}$ (J) & $F_{\rm cross}$ (N) \\
\midrule
Tabletop (1\,L, 10\,T) & 10 & $10^{-3}$ & $3.98 \times 10^4$ &
  $2.2 \times 10^{-10}$ \\
1\,m shell (10\,T) & 10 & 4.19 & $1.67 \times 10^8$ &
  $9.3 \times 10^{-7}$ \\
5\,m shell (10\,T) & 10 & 524 & $2.08 \times 10^{10}$ &
  $1.2 \times 10^{-4}$ \\
5\,m shell (50\,T) & 50 & 524 & $5.21 \times 10^{11}$ &
  $2.9 \times 10^{-3}$ \\
10\,m shell (50\,T) & 50 & 4189 & $4.17 \times 10^{12}$ &
  $2.3 \times 10^{-2}$ \\
\bottomrule
\end{tabular}
\end{center}

\subsection{Field Strength Scaling}

Since $u_{\rm EM} \propto B^2$:
\begin{equation}
\frac{F(B_2)}{F(B_1)} = \left(\frac{B_2}{B_1}\right)^2.
\end{equation}

Going from 10\,T to 50\,T: factor 25.

Going from 10\,T to 100\,T (pulsed): factor 100.

\begin{tcolorbox}[colback=blue!5!white, colframe=blue!50!black,
  title=Volume Scaling Verdict]
Volume scaling alone is insufficient. Even a 10\,m radius shell at 50\,T
produces only 23\,mN. To reach 1\,N requires $V \sim 2 \times 10^5$\,m$^3$
(a sphere of radius 36\,m) at 50\,T --- an impractical superconducting
structure.

\textbf{Scaling summary:}
\begin{itemize}
\item $V$: factor $\sim 5 \times 10^5$ from 1\,L to 5\,m shell
\item $B^2$: factor 25 from 10\,T to 50\,T
\item Combined: factor $\sim 10^7$
\item Total: $10^7 \times 2.2 \times 10^{-10} = 2.2 \times 10^{-3}$\,N
\end{itemize}
Still three orders of magnitude below 1\,N.
\end{tcolorbox}


%=============================================================================
\section{Combined Amplification Analysis}\label{sec:combined}
%=============================================================================

\subsection{The Three Mechanisms Combined}

The three mechanisms are multiplicative (to leading order):
\begin{equation}\label{eq:combined}
\boxed{
F_{\rm total} = F_{\rm static}(B_0, V_0)
  \times \underbrace{\left(\frac{B}{B_0}\right)^2 \frac{V}{V_0}}_{\text{volume/field}}
  \times \underbrace{G_{\rm par}(\lambda, Q, m)}_{\text{parametric}}.
}
\end{equation}

Note the Q-factor is subsumed: it affects $U_{\rm stored}$ which is
already in $F_{\rm static}$.

\subsection{Parameter Space Map}

Using $F_0 = 2.2 \times 10^{-10}$\,N (baseline: 10\,T, 1\,L):

\begin{center}
\begin{tabular}{lccccc}
\toprule
Scenario & $B$ (T) & $V$ (m$^3$) & $G_{\rm par}$ & $F$ (N) & Viable? \\
\midrule
Static tabletop & 10 & $10^{-3}$ & 1 & $2.2 \times 10^{-10}$ & No \\
Static large shell & 50 & 524 & 1 & $2.9 \times 10^{-3}$ & No \\
Parametric tabletop & 10 & $10^{-3}$ & $10^4$ & $2.2 \times 10^{-6}$ & No \\
Parametric 5\,m/50\,T & 50 & 524 & $10^4$ & 29 & \textbf{Yes} \\
Parametric 5\,m/50\,T & 50 & 524 & $10^3$ & 2.9 & \textbf{Yes} \\
Parametric 5\,m/50\,T & 50 & 524 & $10^2$ & 0.29 & Marginal \\
Parametric 10\,m/50\,T & 50 & 4189 & $10^2$ & 2.3 & \textbf{Yes} \\
\bottomrule
\end{tabular}
\end{center}

\subsection{Minimum Parametric Gain for Useful Thrust}

For $F > 1$\,N with a 5\,m shell at 50\,T ($F_{\rm static} = 2.9 \times 10^{-3}$\,N):
\begin{equation}
G_{\rm par,min} = \frac{1}{2.9 \times 10^{-3}} = 345.
\end{equation}

For a 10\,m shell at 50\,T ($F_{\rm static} = 2.3 \times 10^{-2}$\,N):
\begin{equation}
G_{\rm par,min} = \frac{1}{2.3 \times 10^{-2}} = 43.
\end{equation}

\subsection{What $|\lambda - 1|$ Is Required?}

The parametric gain $G_{\rm par}$ depends on $|\lambda - 1|$ through the
ratio of the driven $\psi$-amplitude to the static one. For a rough
scaling:
\begin{equation}
G_{\rm par} \sim \sqrt{\frac{|\lambda - 1|}{|\lambda - 1|_{\rm min}}}
  \times f(Q_\psi),
\end{equation}
where $Q_\psi = \Omega_\psi/(2\gamma_\psi)$ is the $\psi$-mode quality
factor and $f(Q_\psi) \sim \sqrt{Q_\psi}$ near threshold.

For $G_{\rm par} = 345$ and $|\lambda - 1|_{\rm min} = 10^{-14}$:
\begin{equation}
|\lambda - 1| \gtrsim (345)^2 \times 10^{-14}
  \sim 1.2 \times 10^{-9}.
\end{equation}

This is well within the accidental bound. For $G_{\rm par} = 43$:
\begin{equation}
|\lambda - 1| \gtrsim (43)^2 \times 10^{-14}
  \sim 1.8 \times 10^{-11}.
\end{equation}


%=============================================================================
\section{Deep-Space Enhancement: MOND Regime}\label{sec:MOND}
%=============================================================================

\subsection{The MOND Connection}

In DFD, the MOND-like acceleration $a_0 = 2\sqrt{\alpha}\,cH_0$ modifies
the effective gravitational gradient at very low accelerations. For a
device far from any massive body, the ``external field'' is not Earth's $g$
but the background cosmological gradient.

The cross-term force (\ref{eq:force-ka}) becomes:
\begin{equation}
F_{\rm cross}^{(\rm deep\,space)} = \ka\,\frac{U_{\rm EM}}{c^2}\,a_{\rm eff},
\end{equation}
where $a_{\rm eff}$ is the effective gravitational acceleration at the
device's location.

\subsection{Galactic Outskirts}

At galactocentric radii where $a_N < a_0$, the DFD effective acceleration
is enhanced:
\begin{equation}
a_{\rm eff} = \sqrt{a_N \cdot a_0} > a_N.
\end{equation}

For a device in the outer Milky Way at $r = 50$\,kpc:
$a_N \sim 10^{-11}$\,m/s$^2$, $a_0 \sim 1.2 \times 10^{-10}$\,m/s$^2$:
\begin{equation}
a_{\rm eff} = \sqrt{10^{-11} \times 1.2 \times 10^{-10}}
  = \sqrt{1.2 \times 10^{-21}} = 1.1 \times 10^{-10.5}\,\text{m/s}^2
  \approx 3.5 \times 10^{-11}\,\text{m/s}^2.
\end{equation}

This is \textit{less} than Earth's $g = 9.8$\,m/s$^2$, so deep space
\textit{reduces} the cross-term force. The MOND enhancement applies to
the interpolation function but does not exceed the Newtonian value
of $g$ on Earth's surface.

\subsection{Near a Dense Object}

The cross-term force is maximised in strong gravitational fields:
\begin{center}
\begin{tabular}{lcc}
\toprule
Location & $g$ (m/s$^2$) & Enhancement vs Earth \\
\midrule
Earth surface & 9.8 & 1 \\
Jupiter surface & 24.8 & 2.5 \\
White dwarf surface & $\sim 10^6$ & $10^5$ \\
Neutron star surface & $\sim 10^{12}$ & $10^{11}$ \\
\bottomrule
\end{tabular}
\end{center}

Near a neutron star surface, the static force for a 5\,m/50\,T shell would be:
$F \sim 2.9 \times 10^{-3} \times 10^{11} = 2.9 \times 10^8$\,N.
But this is academic --- the device would be destroyed.

\begin{tcolorbox}[colback=red!5!white, colframe=red!50!black,
  title=MOND Regime Verdict]
The deep-space MOND regime does \textbf{not} help. The MOND enhancement
increases the effective gravity at low accelerations but never exceeds
Earth's surface gravity. The cross-term force is maximised in strong-field
environments (near compact objects), not in the MOND regime.

For practical propulsion in free space, the relevant $g$ is the local
gravitational acceleration, which is typically $\ll 9.8$\,m/s$^2$.
\end{tcolorbox}


%=============================================================================
\section{Impossibility Theorem (Without $\lambda \neq 1$)}\label{sec:impossibility}
%=============================================================================

\begin{theorem}[Static Cross-Term Force Bound]\label{thm:bound}
For any static electromagnetic configuration with total energy
$U_{\rm EM}$ in a gravitational field with local acceleration $g$,
the cross-term force is bounded by:
\begin{equation}
\boxed{F_{\rm cross} \leq \ka\,\frac{U_{\rm EM}}{c^2}\,g
  = \frac{3}{8\alpha}\,\frac{U_{\rm EM}}{c^2}\,g.}
\end{equation}
This is exactly the gravitational force on a mass $m_{\rm eff} = \ka\,U_{\rm EM}/c^2$.
\end{theorem}

\begin{proof}
The proof follows from the linearity of $\delta\psi_{\rm EM}$ in the
weak-field regime ($|\delta\psi| \ll 1$). The field equation
$\nabla^2(\delta\psi) = -(8\pi G/c^4)\ka\,u_{\rm EM}$ is linear in $u_{\rm EM}$.
The cross-term force integral is linear in $\delta\psi$, hence linear in
$U_{\rm EM}$. The coefficient $\ka$ is the maximum coupling (achieved when
$\eta > \eta_c$). \qed
\end{proof}

\paragraph{Numerical ceiling (static, Earth).}
To reach $F = 1$\,N requires:
\begin{equation}
U_{\rm EM} = \frac{c^2}{\ka\,g} = \frac{9 \times 10^{16}}{51.4 \times 9.8}
  = 1.79 \times 10^{14}\,\text{J} = 179\,\text{TJ}.
\end{equation}

This is roughly the energy content of 4300 tonnes of TNT, or a
$\sim 100$\,km$^3$ volume at 10\,T. The static channel is
\textit{fundamentally incapable} of producing useful thrust with
realistic EM energies.

%=============================================================================
\section{Summary and Critical Unknowns}\label{sec:summary}
%=============================================================================

\subsection{Force Summary Table}

\begin{center}
\begin{tabular}{llc}
\toprule
Mechanism & Force Expression & Comment \\
\midrule
Static cross-term & $\ka(U_{\rm EM}/c^2)g$ & Baseline, $\sim 10^{-10}$\,N \\
Q-factor & Same as static & No additional amplification \\
Volume scaling & $\propto V$, $\propto B^2$ & Factor $\sim 10^7$ max \\
Parametric ($|\lambda-1| \neq 0$) & $G_{\rm par} \times$ static & $G_{\rm par}$ up to $\sim 10^{13}$ \\
MOND deep space & Replace $g \to a_{\rm eff}$ & Reduces force ($a_{\rm eff} < g$) \\
\bottomrule
\end{tabular}
\end{center}

\subsection{Paths to $F > 1$\,N}

\begin{enumerate}
\item \textbf{Static only:} Requires $U_{\rm EM} > 179$\,TJ.
  \textbf{Impossible} with foreseeable technology.

\item \textbf{Volume + field scaling only ($G_{\rm par} = 1$):}
  Best case 10\,m/50\,T: $F = 23$\,mN. Need 44$\times$ more.
  A 67\,m radius shell at 50\,T: theoretically possible but
  engineering-impossible.

\item \textbf{Parametric + volume + field:}
  5\,m shell, 50\,T, $G_{\rm par} \geq 345$: \textbf{achievable if}
  $|\lambda - 1| \gtrsim 10^{-9}$.

  10\,m shell, 50\,T, $G_{\rm par} \geq 43$: \textbf{achievable if}
  $|\lambda - 1| \gtrsim 2 \times 10^{-11}$.
\end{enumerate}

\subsection{Critical Unknowns}

\begin{enumerate}
\item \textbf{Is $\lambda \neq 1$?}
  This is the single most important unknown. The accidental bound says
  $|\lambda - 1| \lesssim 3 \times 10^{-5}$. If $\lambda = 1$ exactly,
  parametric amplification is impossible and the cross-term force is
  limited to the static value --- forever negligible.

\item \textbf{What is the $\psi$-mode frequency $\Omega_\psi$?}
  The parametric channel requires $2\omega = 2\Omega_\psi$.
  If $\Omega_\psi$ is extremely high or extremely low, it may be
  impractical to match with EM driving.

\item \textbf{What is the $\psi$-mode damping rate $\gamma_\psi$?}
  The parametric gain scales inversely with $\gamma_\psi$ (through
  $Q_\psi = \Omega_\psi/(2\gamma_\psi)$).

\item \textbf{What is the nonlinear saturation physics?}
  Our estimate $G_{\rm par} \sim \sqrt{|\lambda-1|/|\lambda-1|_{\rm min}}$
  is a rough scaling. The actual saturation depends on mode-mode coupling
  coefficients that are not yet computed.
\end{enumerate}

\subsection{Experimental Strategy}

The path forward is unambiguous:
\begin{enumerate}
\item \textbf{First:} Determine whether $\lambda \neq 1$ via the
  intentional detection protocol (Appendix~R, \S R.6). This is a
  single-afternoon measurement using existing cavity apparatus.
\item \textbf{If $\lambda \neq 1$:} Measure $|\lambda - 1|$ precisely.
  If $|\lambda - 1| > 10^{-11}$, parametric amplification of the
  cross-term force could reach useful levels in large shells.
\item \textbf{If $\lambda = 1$:} The cross-term force is forever
  negligible. Move to Channel~3 (direct $\psi$-wave radiation) or
  Channel~5 (vacuum-loading gradient) for propulsion concepts.
\end{enumerate}

\begin{tcolorbox}[colback=green!5!white, colframe=green!50!black,
  title=Key Result]
\textbf{The cross-term self-acceleration force is:}
\begin{equation}
F = 51.4 \times \frac{U_{\rm EM}}{c^2} \times g \times G_{\rm par}.
\end{equation}

Without parametric amplification ($G_{\rm par} = 1$), useful thrust ($>1$\,N)
requires $>179$\,TJ of stored EM energy --- impossible.

With parametric amplification ($G_{\rm par} \sim 10^3$--$10^4$) in a large
(5--10\,m) high-field (50\,T) superconducting shell, thrust of 1--30\,N is
achievable \textbf{if and only if} $|\lambda - 1| \gtrsim 10^{-11}$.

\textbf{The entire propulsion concept stands or falls on the measurement
of $\lambda$.} A single cavity experiment can resolve this.
\end{tcolorbox}


\end{document}
